% GENERATED FILE. Edit canonical lessons, then run npm run generate:books.
% canonical-source-hash: fnv1a64:8e0b712535b0e9d8
% canonical-lessons: BN-C120-choti, BN-C120-maphlar, BN-C120-har, BN-C120-shiyal, BN-C120-badur

\chapter{Slippers, Scarf, Necklace, Jackal, Bat}
\label{ch:bn-slippers-scarf-necklace-jackal-bat}
\hlchaptermodality{120}

\begin{chapteropening}
\textbf{By the end of this chapter:} \emph{I can say slippers, a scarf, a necklace, a jackal and a bat.}

Complete the last lesson of chapter 120: I can say slippers, a scarf, a necklace, a jackal and a bat.
\end{chapteropening}

\section[\texorpdfstring{\bn{চটি} (chôṭi)}{চটি (chôṭi)}]{\bn{চটি} (chôṭi) — slippers}

\label{lesson:BN-C120-choti}



\begin{quote}
Before the new one: say the Bengali for a dhoti, then the Bengali for a long shirt.
\end{quote}

\subsection*{You'll want to know: \bn{চটি}}
\textbf{\bn{চটি}} — \emph{chôṭi} — \textquotedblleft{}slippers\textquotedblright{}.

\begin{cousinweb}
The origin of \textbf{\bn{চটি}} (\emph{chôṭi}) is not settled: CDIAL has no entry for it, and a search by form and by the glosses \textquotedblleft{}slipper\textquotedblright{} and \textquotedblleft{}shoe\textquotedblright{} finds nothing close.
\end{cousinweb}

\begin{grammarlens}[title={What you've built}]
One more word for this chapter.
\end{grammarlens}

\subsection*{Your turn}
\begin{itemize}
\raggedright
  \item \emph{Say it:} \emph{chôṭi}
  \item \emph{Say it:} \emph{chôṭi}, once more
  \item \emph{Say it:} say \emph{chôṭi}
  \item \emph{Recall:} say the Bengali for a dhoti, then the Bengali for a long shirt, then say \emph{chôṭi} again
\end{itemize}

\begin{tcolorbox}[breakable,colback=teal!4,colframe=teal!35!black,title={Before you move on}]
What does \bn{চটি} mean? (\textquotedblleft{}Slippers\textquotedblright{}.) Say it once more.
\end{tcolorbox}
\section[\texorpdfstring{\bn{মাফলার} (māphlār)}{মাফলার (māphlār)}]{\bn{মাফলার} (māphlār) — a scarf}

\label{lesson:BN-C120-maphlar}



\begin{quote}
Before the new one: say the Bengali for a long shirt, then the Bengali for slippers.
\end{quote}

\subsection*{You'll want to know: \bn{মাফলার}}
\textbf{\bn{মাফলার}} — \emph{māphlār} — \textquotedblleft{}a scarf\textquotedblright{}.

\begin{cousinweb}
\textbf{\bn{মাফলার}} (\emph{māphlār}) is a loan from English \emph{muffler}.
\end{cousinweb}

\begin{grammarlens}[title={What you've built}]
One more word for this chapter.
\end{grammarlens}

\subsection*{Your turn}
\begin{itemize}
\raggedright
  \item \emph{Say it:} \emph{māphlār}
  \item \emph{Say it:} \emph{māphlār}, once more
  \item \emph{Say it:} say \emph{māphlār}
  \item \emph{Recall:} say the Bengali for a long shirt, then the Bengali for slippers, then say \emph{māphlār} again
\end{itemize}

\begin{tcolorbox}[breakable,colback=teal!4,colframe=teal!35!black,title={Before you move on}]
What does \bn{মাফলার} mean? (\textquotedblleft{}A scarf\textquotedblright{}.) Say it once more.
\end{tcolorbox}
\section[\texorpdfstring{\bn{হার} (hār)}{হার (hār)}]{\bn{হার} (hār) — a necklace}

\label{lesson:BN-C120-har}



\begin{quote}
Before the new one: say the Bengali for slippers, then the Bengali for a scarf.
\end{quote}

\subsection*{You'll want to know: \bn{হার}}
\textbf{\bn{হার}} — \emph{hār} — \textquotedblleft{}a necklace\textquotedblright{}.

\begin{cousinweb}
\textbf{\bn{হার}} (\emph{hār}) is inherited from Sanskrit \emph{hāra}, \textquotedblleft{}necklace\textquotedblright{}, and has come down with almost no change (CDIAL 14059).
\end{cousinweb}

\begin{grammarlens}[title={What you've built}]
One more word for this chapter.
\end{grammarlens}

\subsection*{Your turn}
\begin{itemize}
\raggedright
  \item \emph{Say it:} \emph{hār}
  \item \emph{Say it:} \emph{hār}, once more
  \item \emph{Say it:} say \emph{hār}
  \item \emph{Recall:} say the Bengali for slippers, then the Bengali for a scarf, then say \emph{hār} again
\end{itemize}

\begin{tcolorbox}[breakable,colback=teal!4,colframe=teal!35!black,title={Before you move on}]
What does \bn{হার} mean? (\textquotedblleft{}A necklace\textquotedblright{}.) Say it once more.
\end{tcolorbox}
\section[\texorpdfstring{\bn{শিয়াল} (shiyāl)}{শিয়াল (shiyāl)}]{\bn{শিয়াল} (shiyāl) — a jackal}

\label{lesson:BN-C120-shiyal}



\begin{quote}
Before the new one: say the Bengali for a scarf, then the Bengali for a necklace.
\end{quote}

\subsection*{You'll want to know: \bn{শিয়াল}}
\textbf{\bn{শিয়াল}} — \emph{shiyāl} — \textquotedblleft{}a jackal\textquotedblright{}.

\begin{cousinweb}
\textbf{\bn{শিয়াল}} (\emph{shiyāl}) is a tadbhava, a word inherited from Sanskrit and reshaped by sound change. It comes from \emph{śṛgāla}, \textquotedblleft{}jackal\textquotedblright{} (CDIAL 12578).
\end{cousinweb}

\begin{grammarlens}[title={What you've built}]
One more word for this chapter.
\end{grammarlens}

\subsection*{Your turn}
\begin{itemize}
\raggedright
  \item \emph{Say it:} \emph{shiyāl}
  \item \emph{Say it:} \emph{shiyāl}, once more
  \item \emph{Say it:} say \emph{shiyāl}
  \item \emph{Recall:} say the Bengali for a scarf, then the Bengali for a necklace, then say \emph{shiyāl} again
\end{itemize}

\begin{tcolorbox}[breakable,colback=teal!4,colframe=teal!35!black,title={Before you move on}]
What does \bn{শিয়াল} mean? (\textquotedblleft{}A jackal\textquotedblright{}.) Say it once more.
\end{tcolorbox}
\section[\texorpdfstring{\bn{বাদুড়} (bādur)}{বাদুড় (bādur)}]{\bn{বাদুড়} (bādur) — a bat}

\label{lesson:BN-C120-badur}



\begin{quote}
Before the new one: say the Bengali for a necklace, then the Bengali for a jackal.
\end{quote}

\subsection*{You'll want to know: \bn{বাদুড়}}
\textbf{\bn{বাদুড়}} — \emph{bādur} — \textquotedblleft{}a bat\textquotedblright{}.

\begin{cousinweb}
\textbf{\bn{বাদুড়}} (\emph{bādur}) is a tadbhava, a word inherited and reshaped by sound change. It comes from \emph{*vālduḍa}, \textquotedblleft{}bat\textquotedblright{}, a form CDIAL reconstructs from its descendants (CDIAL 11584). It stands beside Sanskrit \emph{vālguda}, \textquotedblleft{}bat\textquotedblright{}; CDIAL reports Bloch's view that the word is Dravidian.
\end{cousinweb}

\begin{grammarlens}[title={What you've built}]
One more word for this chapter.
\end{grammarlens}

\subsection*{Your turn}
\begin{itemize}
\raggedright
  \item \emph{Say it:} \emph{bādur}
  \item \emph{Say it:} \emph{bādur}, once more
  \item \emph{Say it:} say \emph{bādur}
  \item \emph{Recall:} say the Bengali for a necklace, then the Bengali for a jackal, then say \emph{bādur} again
\end{itemize}

\begin{tcolorbox}[breakable,colback=teal!4,colframe=teal!35!black,title={Before you move on}]
What does \bn{বাদুড়} mean? (\textquotedblleft{}A bat\textquotedblright{}.) Say it once more.
\end{tcolorbox}
